\documentclass[11pt]{article}
\usepackage{fontspec,amsmath,amssymb,amsthm,hyperref}
\usepackage[margin=1in]{geometry}
\setlength{\emergencystretch}{3em}
\title{Compatibility exports for projection of quadratic transport}
\author{LangevinLean mathematical exposition}
\date{}
\begin{document}\maketitle
\noindent Source module: \texttt{IPLA/Transport.lean}. Compile this proof document with XeLaTeX or LuaLaTeX. The accompanying HTML places the complete Lean declarations beside every proof.


This module exports the shared projected RMS and transport inequalities. The proof lives in \texttt{UnadjustedLangevin/Projection.lean}; IPLA instantiates it with its norm-one parameter projection.


\section{Minkowski on each coupling gives the projected RMS estimate}


For a linear contraction \(T\), any center \(b\), and a coupling \(c\) of \(\mu,\pi\), write \(Tx-b=T(x-y)+(Ty-b)\). Minkowski and \(\|T(x-y)\|\le\|x-y\|\) imply

\[\|Tx-b\|_{L^2(\mu)}\le\operatorname{cost}(c)+\|Ty-b\|_{L^2(\pi)}.\]

Taking the infimum over couplings proves the same inequality with \(W_2(\mu,\pi)\). Finite transport plus a finite equilibrium radius proves finiteness of the projected moment. Only then does the shared real conversion identify its value with \(\sqrt{\int\|Tx-b\|^2d\mu}\). A second-moment bound at equilibrium supplies the concentration radius. No optimal coupling or separate Wasserstein triangle inequality is needed.



\noindent Lean declarations: .

\end{document}
